\section{Introduction to Machine Learning \& Data
Science}\label{introduction-to-machine-learning-data-science}

Welcome to the \textbf{Interactive Machine Learning Curriculum}! Before
diving into code, let's understand what we are trying to achieve. At its
core, \textbf{Artificial Intelligence (AI)} is about teaching computers
to make smart decisions. \textbf{Machine Learning (ML)} is a subset of
AI where, instead of writing strict rules (like \texttt{if-else}
statements), we feed the computer data and let it discover the rules
itself. \textbf{Deep Learning (DL)} takes this a step further by using
layered structures (neural networks) that mimic the human brain.

To program these systems, we follow a simple pipeline: 1.
\textbf{Prepare the Data}: Clean it, scale it, and represent it as
numbers. 2. \textbf{Train a Model}: Let an algorithm search for patterns
in the data. 3. \textbf{Evaluate \& Tune}: Check how well the model
works and adjust its settings. 4. \textbf{Deploy}: Use the model to
predict the future or find hidden structures.

---\n\n\# Chapter 1: Introduction \& Data Manipulation (Pandas \& NumPy)

Welcome to the building blocks of data science! In this unit, we explore
Python's two most essential numerical libraries: \textbf{NumPy}
(Numerical Python) and \textbf{Pandas} (Python Data Analysis Library).
Almost every artificial intelligence or machine learning model you build
starts with loading, manipulating, and cleaning raw data using these
tools.

\subsubsection{High-Level Intuition}\label{high-level-intuition}

\begin{itemize}
\tightlist
\item
  \textbf{NumPy} is like a supercharged calculator. In standard Python,
  if you want to double a list of a million numbers, you have to write a
  loop (\texttt{for\ x\ in\ my\_list}) and process each number one by
  one. This is extremely slow. NumPy introduces the \texttt{ndarray}
  (n-dimensional array), which is a chunk of memory that stores numbers
  close together. It allows us to perform mathematical operations on the
  entire array all at once - a process called \textbf{vectorization}.
\item
  \textbf{Pandas} is like an interactive, programmatic Microsoft Excel
  spreadsheet. It is built directly on top of NumPy. It takes NumPy
  arrays and wraps them in labels, creating \textbf{Series} (1D labeled
  columns) and \textbf{DataFrames} (2D tables with labeled rows and
  columns). With Pandas, you can quickly filter rows, calculate averages
  of groups, and slice datasets.
\end{itemize}

\subsubsection{Core Concepts \& Mathematical
Intuition}\label{core-concepts-mathematical-intuition}

\begin{enumerate}
\def\labelenumi{\arabic{enumi}.}
\tightlist
\item
  \textbf{Vectorization}: Instead of processing items sequentially,
  calculations are compiled in C and run in parallel on CPU registers.
  For example, multiplying an array by 2 (\texttt{arr\ *\ 2}) takes one
  fast CPU instruction.
\item
  \textbf{Data Slicing \& Selection}:

  \begin{itemize}
  \tightlist
  \item
    \texttt{iloc} is integer-based selection
    (e.g.~\texttt{df.iloc{[}0,\ 1{]}} selects the row at position 0 and
    column at position 1).
  \item
    \texttt{loc} is label-based selection
    (e.g.~\texttt{df.loc{[}df{[}\textquotesingle{}Age\textquotesingle{}{]}\ \textgreater{}\ 30,\ \textquotesingle{}Name\textquotesingle{}{]}}
    selects the `Name' column for all rows where `Age' is greater than
    30).
  \end{itemize}
\item
  \textbf{Groupby and Aggregate}: Imagine you have a list of sales
  transactions by product. A
  \texttt{groupby(\textquotesingle{}Product\textquotesingle{})}
  operation groups all matching records together, and an aggregation
  like \texttt{.sum()} computes the total sum of sales for each unique
  product group.
\end{enumerate}

\subsubsection{Programming Guide}\label{programming-guide}

```python import numpy as np import pandas as pd

\section{Creating a NumPy array}\label{creating-a-numpy-array}

arr = np.array({[}1, 2, 3, 4, 5{]}) double\_arr = arr * 2 \# Vectorized
doubling

\section{Creating a Pandas DataFrame from a
dictionary}\label{creating-a-pandas-dataframe-from-a-dictionary}

data = \{`Name': {[}`Alice', `Bob', `Charlie'{]}, `Age': {[}25, 30,
35{]}\} df = pd.DataFrame(data)

\section{Filtering and Selecting}\label{filtering-and-selecting}

age\_col = df{[}`Age'{]} \# Select column filtered\_df =
df.loc{[}df{[}`Age'{]} \textgreater{} 28{]} \# Select rows meeting
criteria ```\n\n---\n\n\# Chapter 2: Data Cleaning \& Preprocessing

Machine learning algorithms are mathematical equations. They cannot read
raw text, they get confused by missing values (\texttt{NaN}), and they
can easily be skewed if one column contains tiny numbers (like age)
while another contains huge numbers (like salary). \textbf{Data
preprocessing} is the process of washing, slicing, and preparing your
raw data before feeding it to a model.

\subsubsection{High-Level Intuition}\label{high-level-intuition-1}

Imagine preparing raw vegetables for cooking: you wash off the dirt,
peel off the bad spots, and chop everything into uniform sizes. In data
cleaning: * \textbf{Missing values} are the rotten spots - we must
either throw them away (drop rows) or replace them with a sensible value
like the average (impute them). * \textbf{Categorical variables} (like
`City' or `Size') are text tags - we must translate them into numbers
(encode them) so the computer can compute with them. * \textbf{Feature
Scaling} is making sure all features share a standard range - so columns
with larger numerical scales don't dominate the mathematical equations.

\subsubsection{Core Concepts \& Mathematical
Intuition}\label{core-concepts-mathematical-intuition-1}

\begin{enumerate}
\def\labelenumi{\arabic{enumi}.}
\tightlist
\item
  \textbf{Handling Missing Values}:

  \begin{itemize}
  \tightlist
  \item
    Removing rows with \texttt{dropna()} is safe if only a few rows have
    missing data, but it wastes valuable information if used
    excessively.
  \item
    Imputing with the median/mean using \texttt{fillna()} fills in gaps
    without changing the overall statistical properties of the dataset.
  \end{itemize}
\item
  \textbf{Encoding Categorical Columns}:

  \begin{itemize}
  \tightlist
  \item
    \textbf{Ordinal Encoding}: Used when categories have a natural
    sequence (e.g., Low=0, Medium=1, High=2). We map them directly to
    integer series.
  \item
    \textbf{One-Hot Encoding}: Used when categories have no logical
    sequence (e.g., Red, Blue, Green). We create dummy binary columns
    for each category (e.g., \texttt{Color\_Red}, \texttt{Color\_Blue}),
    so the model doesn't assume that Blue is ``greater than'' Red.
  \end{itemize}
\item
  \textbf{Feature Scaling}:

  \begin{itemize}
  \tightlist
  \item
    \textbf{Min-Max Scaling}: Squashes values into a \([0, 1]\) range.
    Formula: \(x_{scaled} = \frac{x - x_{min}}{x_{max} - x_{min}}\).
  \item
    \textbf{Standardization (Z-score)}: Centers data around 0 with a
    standard deviation of 1. Formula: \(z = \frac{x - \mu}{\sigma}\)
    (where \(\mu\) is the mean and \(\sigma\) is the standard
    deviation).
  \end{itemize}
\end{enumerate}

\subsubsection{Programming Guide}\label{programming-guide-1}

```python from sklearn.preprocessing import MinMaxScaler, StandardScaler
import pandas as pd

\section{Filling missing values in
Pandas}\label{filling-missing-values-in-pandas}

df{[}`Score'{]} = df{[}`Score'{]}.fillna(df{[}`Score'{]}.mean())

\section{One-hot encoding categorical
variables}\label{one-hot-encoding-categorical-variables}

df\_encoded = pd.get\_dummies(df, columns={[}`Category'{]})

\section{Scaling features to {[}0, 1{]}}\label{scaling-features-to-0-1}

scaler = MinMaxScaler() scaled\_features =
scaler.fit\_transform(df{[}{[}`Age', `Income'{]}{]}) ```\n\n---\n\n\#
Chapter 3: Data Visualization \& Plotting

Data visualization is key to exploratory data analysis (EDA). Before
building a machine learning model, you must use your own eyes to study
the distribution, trends, correlation, and anomalies in your data.

\subsubsection{High-Level Intuition}\label{high-level-intuition-2}

If you look at a spreadsheet of 100,000 house prices, you won't see
anything but numbers. But if you plot a \textbf{scatter plot} of house
size vs.~price, you will immediately see if the relationship is linear,
curved, or if there are weird outliers (like a tiny house selling for
\$10 million). * \textbf{Matplotlib} is your drawing board - it gives
you low-level controls over every line, tick mark, and text label. *
\textbf{Seaborn} is your high-level artist - it takes DataFrames
directly and creates gorgeous, statistically rich plots automatically.

\subsubsection{Core Concepts}\label{core-concepts}

\begin{enumerate}
\def\labelenumi{\arabic{enumi}.}
\tightlist
\item
  \textbf{Histograms}: Plots the distribution of a single numerical
  column by grouping values into bins, showing you if the data is
  symmetric or skewed.
\item
  \textbf{Scatter Plots with Regression}: Shows the relationship between
  two variables, highlighting how they correlate. Adding a trendline
  highlights the general direction.
\item
  \textbf{Correlation Heatmaps}: A grid showing the correlation
  coefficient (\(r\)) between every pair of features in the dataset. We
  mask the upper diagonal to avoid redundant information.
\end{enumerate}

\subsubsection{Programming Guide}\label{programming-guide-2}

```python import matplotlib.pyplot as plt import seaborn as sns

\section{Basic line plot in
Matplotlib}\label{basic-line-plot-in-matplotlib}

plt.plot(x, y) plt.xlabel(`Time') plt.ylabel(`Value') plt.title(`Line
Trend') plt.show()

\section{Scatter plot with regression line in
Seaborn}\label{scatter-plot-with-regression-line-in-seaborn}

sns.lmplot(x=`Age', y=`Salary', data=df)

\section{Correlation heatmap}\label{correlation-heatmap}

corr = df.corr() sns.heatmap(corr, annot=True, cmap=`coolwarm')
```\n\n---\n\n\# Chapter 4: Regression (Linear \& Logistic)

Regression is the branch of machine learning where we predict values. *
\textbf{Linear Regression} predicts continuous numbers (like house
prices or temperatures). * \textbf{Logistic Regression} predicts
probabilities for binary categories (like whether an email is spam (1)
or not spam (0)).

\subsubsection{High-Level Intuition}\label{high-level-intuition-3}

\begin{itemize}
\tightlist
\item
  \textbf{Linear Regression} is like drawing a straight line through a
  cloud of points on a graph so that the total distance from the points
  to the line is as small as possible.
\item
  \textbf{Logistic Regression} is like wrapping that straight line
  around an S-shaped curve (called a sigmoid function). This curve
  squashes any value between \(-\infty\) and \(+\infty\) into a
  probability between \(0\) and \(1\).
\item
  \textbf{Gradient Descent} is like being blindfolded on a foggy
  mountain and taking steps downward in the direction of the steepest
  slope until you reach the bottom of the valley (the minimum error).
\end{itemize}

\subsubsection{Core Concepts \& Mathematical
Intuition}\label{core-concepts-mathematical-intuition-2}

\begin{enumerate}
\def\labelenumi{\arabic{enumi}.}
\tightlist
\item
  \textbf{Linear Regression Formula}: \(y = w \cdot x + b\). We find the
  weights (\(w\)) and bias (\(b\)) that minimize the Mean Squared Error
  (MSE).
\item
  \textbf{Logistic Regression Sigmoid}:
  \(P(y=1) = \frac{1}{1 + e^{-z}}\) where \(z = w \cdot x + b\). If
  \(P \ge 0.5\), we classify as class 1, otherwise class 0.
\item
  \textbf{Regularization}: Prevents the model from overfitting
  (memorizing training data):

  \begin{itemize}
  \tightlist
  \item
    \textbf{L1 (Lasso)}: Adds a penalty proportional to the absolute
    values of the weights (\(\sum |w_i|\)). It shrinks unimportant
    weights to exactly 0, acting as a feature selector.
  \item
    \textbf{L2 (Ridge)}: Adds a penalty proportional to the squared
    values of the weights (\(\sum w_i^2\)). It shrinks weights close to
    0, reducing model complexity.
  \end{itemize}
\end{enumerate}

\subsubsection{Programming Guide}\label{programming-guide-3}

```python from sklearn.linear\_model import LinearRegression,
LogisticRegression, Lasso, Ridge

\section{Fitting linear regression}\label{fitting-linear-regression}

lr = LinearRegression().fit(X\_train, y\_train) y\_pred =
lr.predict(X\_test)

\section{Fitting L2 Regularized Ridge
Regression}\label{fitting-l2-regularized-ridge-regression}

ridge = Ridge(alpha=1.0).fit(X\_train, y\_train) ```\n\n---\n\n\#
Chapter 5: Decision Trees \& Random Forests

Decision Trees and Random Forests are intuitive, rule-based algorithms
used for classification and regression tasks.

\subsubsection{High-Level Intuition}\label{high-level-intuition-4}

\begin{itemize}
\tightlist
\item
  \textbf{Decision Tree} is like playing a game of ``20 Questions.'' You
  start at the root and ask a simple question (e.g.~``Is age
  \textgreater{} 30?''). Depending on the answer, you go left or right
  and ask another question, until you reach a leaf node (the
  prediction).
\item
  \textbf{Random Forest} is the wisdom of the crowd. Instead of trusting
  a single decision tree (which might be biased or memorize details), we
  train a forest of 100 trees on random subsets of the data and
  features. To get a final prediction, we let all 100 trees vote and
  take the majority decision.
\end{itemize}

\subsubsection{Core Concepts}\label{core-concepts-1}

\begin{enumerate}
\def\labelenumi{\arabic{enumi}.}
\tightlist
\item
  \textbf{Information Gain \& Gini Impurity}: Decision trees choose
  splits that maximize ``purity'' (keeping similar classes together in
  the resulting branches).
\item
  \textbf{Overfitting}: Trees left to grow without limits will create
  rules for every single sample, resulting in poor general performance.
  We prevent this by setting \texttt{max\_depth} to prune the tree.
\item
  \textbf{Feature Importance}: Trees naturally track how often a feature
  is used to split the data. Features that make clean splits near the
  top of the tree are marked as highly important.
\end{enumerate}

\subsubsection{Programming Guide}\label{programming-guide-4}

```python from sklearn.tree import DecisionTreeClassifier from
sklearn.ensemble import RandomForestClassifier

\section{Training a constrained Decision
Tree}\label{training-a-constrained-decision-tree}

dt = DecisionTreeClassifier(max\_depth=3, random\_state=42)
dt.fit(X\_train, y\_train)

\section{Training a Random Forest}\label{training-a-random-forest}

rf = RandomForestClassifier(n\_estimators=100, max\_depth=5,
random\_state=42) rf.fit(X\_train, y\_train) importances =
rf.feature\_importances\_ ```\n\n---\n\n\# Chapter 6: Support Vector
Machines \& Naive Bayes

In this unit, we explore two highly effective classification algorithms,
particularly powerful for text mining and boundary separation.

\subsubsection{High-Level Intuition}\label{high-level-intuition-5}

\begin{itemize}
\tightlist
\item
  \textbf{Naive Bayes} is based on word frequencies. To classify whether
  an email is spam, it looks at the probability of each word appearing
  in spam emails vs.~non-spam emails. It is ``naive'' because it assumes
  words don't influence each other (e.g.~it treats ``New'' and ``York''
  as completely independent), but it is incredibly fast and works well
  for text.
\item
  \textbf{Support Vector Machine (SVM)} is about drawing borders.
  Imagine black and white marbles scattered on a table. SVM draws a
  straight line that separates them while maximizing the empty space
  (the margin) between the line and the closest marbles (support
  vectors).
\end{itemize}

\subsubsection{Core Concepts \& Mathematical
Intuition}\label{core-concepts-mathematical-intuition-3}

\begin{enumerate}
\def\labelenumi{\arabic{enumi}.}
\tightlist
\item
  \textbf{Bayes' Theorem}: Calculates
  \(P(Class | Word) = \frac{P(Word | Class) \cdot P(Class)}{P(Word)}\).
\item
  \textbf{TF-IDF}: Term Frequency-Inverse Document Frequency. Scores a
  word's importance by multiplying how often it appears in a document by
  how rare it is across all documents.
\item
  \textbf{Kernel Trick}: When marbles are not separable by a straight
  line, SVM uses a kernel (like the RBF kernel) to map data into a
  higher dimension (like lifting them into 3D space) where a flat sheet
  can easily cut between them.
\end{enumerate}

\subsubsection{Programming Guide}\label{programming-guide-5}

```python from sklearn.feature\_extraction.text import TfidfVectorizer
from sklearn.naive\_bayes import MultinomialNB from sklearn.svm import
SVC

\section{Converting text corpus to TF-IDF
features}\label{converting-text-corpus-to-tf-idf-features}

vec = TfidfVectorizer() X\_tfidf = vec.fit\_transform(text\_corpus)

\section{Fitting Naive Bayes}\label{fitting-naive-bayes}

nb = MultinomialNB().fit(X\_tfidf, y)

\section{Fitting SVM with non-linear RBF
kernel}\label{fitting-svm-with-non-linear-rbf-kernel}

svm = SVC(kernel=`rbf', C=1.0).fit(X\_train, y\_train) ```\n\n---\n\n\#
Chapter 7: Unsupervised Learning \& Clustering

Unsupervised learning is used when we have inputs but no pre-existing
answers (labels). We want the computer to discover hidden structures on
its own.

\subsubsection{High-Level Intuition}\label{high-level-intuition-6}

\begin{itemize}
\tightlist
\item
  \textbf{K-Means Clustering} is like sorting scattered objects into
  \(K\) separate piles. We randomly drop \(K\) flags (centroids), assign
  objects to the nearest flag, move each flag to the center of its pile,
  and repeat until the piles stop moving.
\item
  \textbf{PCA (Principal Component Analysis)} is like photographing a 3D
  object. A camera drops the depth dimension (3D to 2D) but chooses an
  angle that captures the object's width and height as clearly as
  possible. PCA finds the directions of highest spread (variance) to
  compress variables.
\end{itemize}

\subsubsection{Core Concepts}\label{core-concepts-2}

\begin{enumerate}
\def\labelenumi{\arabic{enumi}.}
\tightlist
\item
  \textbf{Centroids and Inertia}: Inertia measures the sum of squared
  distances from data points to their assigned cluster centroid. We want
  this distance to be small.
\item
  \textbf{Elbow Method}: We plot inertia for different values of \(K\).
  The ``elbow'' point is where adding more clusters yields diminishing
  returns in reducing inertia.
\item
  \textbf{Dimensionality Reduction}: Shrinking the number of variables
  to simplify models, speed up training, and visualize data in 2D or 3D.
\end{enumerate}

\subsubsection{Programming Guide}\label{programming-guide-6}

```python from sklearn.cluster import KMeans from sklearn.decomposition
import PCA

\section{Clustering with K-Means}\label{clustering-with-k-means}

kmeans = KMeans(n\_clusters=3, random\_state=42).fit(X) labels =
kmeans.labels\_

\section{Compressing data using PCA}\label{compressing-data-using-pca}

pca = PCA(n\_components=2) X\_reduced = pca.fit\_transform(X)
```\n\n---\n\n\# Chapter 8: Feature Engineering \& Selection

Feature engineering is the process of using domain knowledge to create
new features (clues) that make machine learning algorithms work better.

\subsubsection{High-Level Intuition}\label{high-level-intuition-7}

\begin{itemize}
\tightlist
\item
  \textbf{Feature Engineering} is like adding helpful signs on a road.
  For example, if you want a model to predict house heating bills,
  giving it \texttt{Length} and \texttt{Width} separately is okay, but
  multiplying them to calculate \texttt{Area} is a much stronger clue.
\item
  \textbf{Feature Selection} is like cleaning out your backpack before a
  hike. You remove heavy, useless items (noisy columns) so you can walk
  faster and stay focused.
\end{itemize}

\subsubsection{Core Concepts}\label{core-concepts-3}

\begin{enumerate}
\def\labelenumi{\arabic{enumi}.}
\tightlist
\item
  \textbf{Polynomial \& Interaction Features}: Creating combinations of
  variables (e.g.~\(x_1^2\), \(x_2^2\), \(x_1 x_2\)) to capture curved
  relationships.
\item
  \textbf{Cyclical Transforms}: Mapping time-based columns (like
  \texttt{hour} from 1 to 24) onto a circle using \(\sin\) and \(\cos\).
  This teaches the model that hour 23 (11 PM) and hour 0 (midnight) are
  close neighbors.
\item
  \textbf{Recursive Feature Elimination (RFE)}: Fits a model, ranks
  features, discards the least useful one, and repeats until the desired
  number of features is reached.
\end{enumerate}

\subsubsection{Programming Guide}\label{programming-guide-7}

```python from sklearn.preprocessing import PolynomialFeatures from
sklearn.feature\_selection import RFE from sklearn.linear\_model import
LinearRegression import numpy as np

\section{Cyclical mapping}\label{cyclical-mapping}

df{[}`sin\_hour'{]} = np.sin(2 * np.pi * df{[}`hour'{]} / 24)
df{[}`cos\_hour'{]} = np.cos(2 * np.pi * df{[}`hour'{]} / 24)

\section{Applying RFE}\label{applying-rfe}

selector = RFE(estimator=LinearRegression(), n\_features\_to\_select=3)
X\_selected = selector.fit\_transform(X, y) ```\n\n---\n\n\# Chapter 9:
Model Evaluation \& Hyperparameter Tuning

How do we know if our model is actually good? In this unit, we explore
evaluation strategies and optimization techniques.

\subsubsection{High-Level Intuition}\label{high-level-intuition-8}

\begin{itemize}
\tightlist
\item
  \textbf{Cross-Validation} is like giving a student multiple practice
  exams. Instead of testing on a single split, we partition data into
  \(K\) blocks. We train on \(K-1\) blocks and test on the remaining
  block, repeating this \(K\) times so every data point is tested once.
\item
  \textbf{Hyperparameter Tuning} is like dialing in the settings of an
  engine. We don't tune weights (parameters); we tune parameters that
  control the model structure (hyperparameters), like tree depth.
\end{itemize}

\subsubsection{Core Concepts}\label{core-concepts-4}

\begin{enumerate}
\def\labelenumi{\arabic{enumi}.}
\tightlist
\item
  \textbf{K-Fold CV}: Reduces evaluation bias by averaging performance
  across \(K\) independent runs.
\item
  \textbf{Grid Search vs.~Random Search}:

  \begin{itemize}
  \tightlist
  \item
    \textbf{Grid Search} checks every possible combination of settings
    from a list. It is thorough but slow.
  \item
    \textbf{Random Search} randomly picks settings from ranges. It is
    fast and often finds optimal configurations in a fraction of the
    time.
  \end{itemize}
\item
  \textbf{Classification Threshold Tuning}: In binary classification,
  changing the probability threshold (default 0.5) lets you balance
  precision (avoiding false alarms) and recall (avoiding missed
  detections).
\end{enumerate}

\subsubsection{Programming Guide}\label{programming-guide-8}

```python from sklearn.model\_selection import cross\_val\_score,
GridSearchCV from sklearn.ensemble import RandomForestClassifier

\section{Evaluating with 5-Fold
Cross-Validation}\label{evaluating-with-5-fold-cross-validation}

scores = cross\_val\_score(RandomForestClassifier(), X, y, cv=5)

\section{Tuning using Grid Search}\label{tuning-using-grid-search}

grid = GridSearchCV( estimator=RandomForestClassifier(),
param\_grid=\{`max\_depth': {[}3, 5{]}, `n\_estimators': {[}10, 50{]}\},
cv=3 ).fit(X, y) best\_model = grid.best\_estimator\_ ```\n\n---\n\n\#
Chapter 10: Introduction to Deep Learning

Deep Learning uses multilayered neural networks to solve complex
regression and classification tasks.

\subsubsection{High-Level Intuition}\label{high-level-intuition-9}

Deep learning is modeled after biological brains. * A \textbf{neuron
(Perceptron)} takes multiple inputs, multiplies them by weights
(importance), sums them up, and decides whether to send a signal forward
using an activation function. * A \textbf{Multi-Layer Perceptron (MLP)}
stacks these neurons in layers. * \textbf{Backpropagation} is the way
neural networks learn: the model makes a guess (forward pass),
calculates how wrong it was (loss), and propagates the error backward
through the layers to adjust weights.

\subsubsection{Core Concepts \& Math}\label{core-concepts-math}

\begin{enumerate}
\def\labelenumi{\arabic{enumi}.}
\tightlist
\item
  \textbf{Perceptron Activation}: \(y = f(\sum w_i x_i + b)\) where
  \(f\) is an activation function (like ReLU (\(max(0,x)\)) or Sigmoid)
  that introduces non-linearity.
\item
  \textbf{Backpropagation \& Gradient Descent}: Updates weights using
  the chain rule:
  \(w \leftarrow w - \alpha \frac{\partial L}{\partial w}\) (where
  \(\alpha\) is the learning rate).
\item
  \textbf{Multi-Layer Architectures}: Stacking hidden layers allows the
  network to learn hierarchical features (e.g.~edges in the first layer,
  shapes in the second, faces in the third).
\end{enumerate}

\subsubsection{Programming Guide}\label{programming-guide-9}

```python from sklearn.neural\_network import MLPClassifier

\section{Building a basic 2-layer Neural
Network}\label{building-a-basic-2-layer-neural-network}

mlp = MLPClassifier( hidden\_layer\_sizes=(64, 32), activation=`relu',
solver=`adam', max\_iter=1000, random\_state=42 ) mlp.fit(X\_train,
y\_train) ```\n\n---\n\n\# Chapter 11: Time Series Analysis \&
Forecasting

Time series data consists of observations ordered sequentially in time.

\subsubsection{High-Level Intuition}\label{high-level-intuition-10}

Imagine predicting tomorrow's temperature. It is highly likely to be
similar to today's temperature (autocorrelation), and it follows
seasonal cycles. Forecasting requires analyzing trends (long-term
directions) and seasonal patterns. We must ensure the series is
``stationary'' - meaning its average value and variance don't drift over
time - before fitting forecasting models.

\subsubsection{Core Concepts}\label{core-concepts-5}

\begin{enumerate}
\def\labelenumi{\arabic{enumi}.}
\tightlist
\item
  \textbf{Stationarity \& ADF Test}: A stationary series is easier to
  predict. The \textbf{Augmented Dickey-Fuller (ADF)} test checks for
  stationarity. If the p-value is \(< 0.05\), the series is stationary.
\item
  \textbf{Differencing}: Subtracting yesterday's value from today's
  value (\(y'_t = y_t - y_{t-1}\)) to remove linear trends and make the
  series stationary.
\item
  \textbf{ARIMA (Autoregressive Integrated Moving Average)}:

  \begin{itemize}
  \tightlist
  \item
    \textbf{AR (p)}: Predicts future values based on past values.
  \item
    \textbf{I (d)}: Number of differencing steps needed to make the data
    stationary.
  \item
    \textbf{MA (q)}: Predicts future values based on past forecast
    errors.
  \end{itemize}
\end{enumerate}

\subsubsection{Programming Guide}\label{programming-guide-10}

```python from statsmodels.tsa.stattools import adfuller from
statsmodels.tsa.arima.model import ARIMA

\section{Checking stationarity}\label{checking-stationarity}

result = adfuller(df{[}`Sales'{]}) print(``p-value:'', result{[}1{]})

\section{Fitting an ARIMA(1, 1, 1)
model}\label{fitting-an-arima1-1-1-model}

model = ARIMA(df{[}`Sales'{]}, order=(1, 1, 1)).fit() forecast =
model.forecast(steps=7) ```\n\n---\n\n\# Chapter 12: Natural Language
Processing (NLP)

Natural Language Processing (NLP) focuses on teaching computers to read,
write, and understand human language.

\subsubsection{High-Level Intuition}\label{high-level-intuition-11}

To a computer, a word is just a string of characters. We must turn text
into numbers. * First, we clean text by lowercasing, removing
punctuation, and dropping generic words like ``the'' or ``is''
(\textbf{stopwords}). * Then, we tokenize sentences into words. * We
represent words as vectors (\textbf{Embeddings}) where words with
similar meanings (e.g.~``king'' and ``queen'') point in similar
directions in space.

\subsubsection{Core Concepts}\label{core-concepts-6}

\begin{enumerate}
\def\labelenumi{\arabic{enumi}.}
\tightlist
\item
  \textbf{TF-IDF Vectorization}: Represents a document as a vector of
  word importances. Highly unique words get high scores.
\item
  \textbf{Cosine Similarity}: Measures the angle between two text
  vectors. Formula:
  \(\cos(\theta) = \frac{\mathbf{u} \cdot \mathbf{v}}{\|\mathbf{u}\| \|\mathbf{v}\|}\).
  If the vectors point in the same direction, similarity is 1.
\item
  \textbf{Word Embeddings}: High-dimensional representations capturing
  semantic relationships between words.
\end{enumerate}

\subsubsection{Programming Guide}\label{programming-guide-11}

```python from sklearn.feature\_extraction.text import TfidfVectorizer
from sklearn.metrics.pairwise import cosine\_similarity

\section{TF-IDF Vectorization}\label{tf-idf-vectorization}

vec = TfidfVectorizer() tfidf\_matrix = vec.fit\_transform(corpus)

\section{Computing similarity between
documents}\label{computing-similarity-between-documents}

sim = cosine\_similarity(tfidf\_matrix{[}0{]}, tfidf\_matrix{[}1{]})
```\n\n---\n\n\# Chapter 13: Anomaly Detection

Anomaly detection is the process of identifying rare, unexpected
occurrences in data.

\subsubsection{High-Level Intuition}\label{high-level-intuition-12}

Anomaly detection is like looking for fraud in credit card transactions.
* Statistical methods use \textbf{Z-scores} or the \textbf{Interquartile
Range (IQR)} to flag values that lie far from the average. *
\textbf{Isolation Forests} work on a simple idea: if you want to isolate
a point in a crowd, a point that is far away (an outlier) requires very
few splits to isolate, while points in dense clusters require many
splits.

\subsubsection{Core Concepts \& Math}\label{core-concepts-math-1}

\begin{enumerate}
\def\labelenumi{\arabic{enumi}.}
\tightlist
\item
  \textbf{Z-Score Limits}: Outliers are points that lie more than 3
  standard deviations away from the mean (\(|z| > 3\)).
\item
  \textbf{IQR Limits}: Outliers lie outside
  \([Q1 - 1.5 \cdot IQR, Q3 + 1.5 \cdot IQR]\).
\item
  \textbf{Mahalanobis Distance}: Measures distance of a point to a
  distribution, accounting for correlation:
  \(D_M(\mathbf{x}) = \sqrt{(\mathbf{x} - \boldsymbol{\mu})^T \boldsymbol{\Sigma}^{-1} (\mathbf{x} - \boldsymbol{\mu})}\)
  (where \(\boldsymbol{\Sigma}\) is the covariance matrix).
\end{enumerate}

\subsubsection{Programming Guide}\label{programming-guide-12}

```python from sklearn.ensemble import IsolationForest import numpy as
np

\section{Fitting Isolation Forest}\label{fitting-isolation-forest}

iso = IsolationForest(contamination=0.02, random\_state=42) outliers =
iso.fit\_predict(X) \# returns -1 for outliers, 1 for normal
```\n\n---\n\n\# Chapter 14: Ensemble Learning \& Gradient Boosting

Ensemble learning combines predictions from multiple individual models
to create a stronger, more robust model.

\subsubsection{High-Level Intuition}\label{high-level-intuition-13}

\begin{itemize}
\tightlist
\item
  \textbf{Voting Classifiers} let different models (e.g.~Logistic
  Regression, SVM, Decision Tree) vote on the prediction. Hard voting
  uses majority rule, while soft voting averages the predicted
  probabilities.
\item
  \textbf{Boosting} trains models sequentially. The first model makes
  predictions, the second model is trained specifically on the errors
  (residuals) of the first, the third on the errors of the second, and
  so on.
\end{itemize}

\subsubsection{Core Concepts}\label{core-concepts-7}

\begin{enumerate}
\def\labelenumi{\arabic{enumi}.}
\tightlist
\item
  \textbf{Soft vs.~Hard Voting}: Soft voting is generally preferred
  because it weighs the confidence of each model's prediction.
\item
  \textbf{Boosting}: Reduces bias by sequentially focusing on
  hard-to-predict samples.
\item
  \textbf{Gradient Boosting}: An additive model where new trees are
  trained to fit the gradient of the loss function.
\end{enumerate}

\subsubsection{Programming Guide}\label{programming-guide-13}

```python from sklearn.ensemble import VotingClassifier,
GradientBoostingClassifier from sklearn.linear\_model import
LogisticRegression from sklearn.svm import SVC

\section{Creating a Soft Voting
Classifier}\label{creating-a-soft-voting-classifier}

voting = VotingClassifier( estimators={[}(`lr', LogisticRegression()),
(`svm', SVC(probability=True)){]}, voting=`soft' ) voting.fit(X\_train,
y\_train)

\section{Gradient Boosting
Classifier}\label{gradient-boosting-classifier}

gb = GradientBoostingClassifier(n\_estimators=100, learning\_rate=0.1)
gb.fit(X\_train, y\_train) ```\n\n---\n\n\# Chapter 15: Dimensionality
Reduction \& Manifold Learning

This unit covers methods for projecting high-dimensional data into lower
dimensions.

\subsubsection{High-Level Intuition}\label{high-level-intuition-14}

\begin{itemize}
\tightlist
\item
  \textbf{LDA (Linear Discriminant Analysis)} is like PCA, but it knows
  about classes. It finds the projection that keeps different groups as
  far apart as possible.
\item
  \textbf{t-SNE} is a visualization tool that maps high-dimensional
  points to 2D. It preserves local neighborhoods, making it easy to see
  clusters visually.
\end{itemize}

\subsubsection{Core Concepts}\label{core-concepts-8}

\begin{enumerate}
\def\labelenumi{\arabic{enumi}.}
\tightlist
\item
  \textbf{LDA (Supervised)}: Projects data to maximize between-class
  variance relative to within-class variance.
\item
  \textbf{t-SNE (Non-linear)}: Converts distances between points to
  probabilities and minimizes divergence between high-dimensional and
  low-dimensional representations.
\item
  \textbf{Silhouette Score}: Measures cluster quality (how close points
  are to their own cluster vs.~other clusters).
\end{enumerate}

\subsubsection{Programming Guide}\label{programming-guide-14}

```python from sklearn.discriminant\_analysis import
LinearDiscriminantAnalysis from sklearn.manifold import TSNE from
sklearn.metrics import silhouette\_score

\section{Running LDA}\label{running-lda}

lda = LinearDiscriminantAnalysis(n\_components=1) X\_lda =
lda.fit\_transform(X, y)

\section{Running t-SNE}\label{running-t-sne}

tsne = TSNE(n\_components=2, perplexity=30, random\_state=42) X\_2d =
tsne.fit\_transform(X) ```\n\n---\n\n\# Chapter 16: Recommender Systems

Recommender systems recommend items (like books or movies) to users
based on preferences.

\subsubsection{High-Level Intuition}\label{high-level-intuition-15}

\begin{itemize}
\tightlist
\item
  \textbf{Content-Based}: Recommending items similar to what the user
  liked before (e.g.~recommending space movies to someone who watched
  \emph{Interstellar}).
\item
  \textbf{Collaborative Filtering}: Finding users with similar tastes
  and recommending items they liked.
\item
  \textbf{Matrix Factorization}: Decomposing a large, sparse user-item
  rating table into user-preference and item-attribute factors using
  Stochastic Gradient Descent (SGD).
\end{itemize}

\subsubsection{Core Concepts}\label{core-concepts-9}

\begin{enumerate}
\def\labelenumi{\arabic{enumi}.}
\tightlist
\item
  \textbf{Cosine Similarity}: Used to compare item feature profiles.
\item
  \textbf{Latent Factors}: Hidden attributes (like movie genre weights)
  that explain user ratings.
\item
  \textbf{Stochastic Gradient Descent (SGD)}: Used to fill in missing
  values in the rating grid by optimizing matrix values iteratively.
\end{enumerate}

\subsubsection{Programming Guide}\label{programming-guide-15}

```python from sklearn.metrics.pairwise import cosine\_similarity import
numpy as np

\section{Calculating item-to-item
similarity}\label{calculating-item-to-item-similarity}

item\_features = np.array({[}{[}1, 0, 1{]}, {[}0, 1, 1{]}, {[}1, 0,
0{]}{]}) similarity = cosine\_similarity(item\_features)
```\n\n---\n\n\# Chapter 17: Association Rule Learning

Association Rule Learning discovers interesting relationships between
items in large transaction databases.

\subsubsection{High-Level Intuition}\label{high-level-intuition-16}

This is the ``market basket analysis'' technique (e.g.~discovering that
people who buy diapers also frequently buy beer).

\subsubsection{Core Concepts \& Math}\label{core-concepts-math-2}

\begin{enumerate}
\def\labelenumi{\arabic{enumi}.}
\tightlist
\item
  \textbf{Support}: How popular an itemset is:
  \[\text{Support}(A) = \frac{\text{Transactions containing } A}{\text{Total Transactions}}\]
\item
  \textbf{Confidence}: How reliable a rule is:
  \[\text{Confidence}(A \rightarrow B) = \frac{\text{Support}(A \cup B)}{\text{Support}(A)}\]
\item
  \textbf{Lift}: How much more likely \(B\) is bought given \(A\),
  compared to buying \(B\) on its own:
  \[\text{Lift}(A \rightarrow B) = \frac{\text{Confidence}(A \rightarrow B)}{\text{Support}(B)}\]
\end{enumerate}

\subsubsection{Programming Guide}\label{programming-guide-16}

\texttt{python\ \#\ Extracting\ association\ metrics\ in\ Pandas\ support\_A\ =\ df{[}\textquotesingle{}A\textquotesingle{}{]}.mean()\ support\_A\_B\ =\ (df{[}\textquotesingle{}A\textquotesingle{}{]}\ \&\ df{[}\textquotesingle{}B\textquotesingle{}{]}).mean()\ confidence\ =\ support\_A\_B\ /\ support\_A\ lift\ =\ confidence\ /\ df{[}\textquotesingle{}B\textquotesingle{}{]}.mean()}\n\n---\n\n\#
Chapter 18: Semi-Supervised Learning

Semi-supervised learning uses a small amount of labeled data and a large
amount of unlabeled data.

\subsubsection{High-Level Intuition}\label{high-level-intuition-17}

Labeling data is expensive. If we have a few labeled points and many
unlabeled points, we can construct a graph connecting similar points and
let the labels ``spread'' across the network like dye in water
(\textbf{Label Propagation}).

\subsubsection{Core Concepts}\label{core-concepts-10}

\begin{enumerate}
\def\labelenumi{\arabic{enumi}.}
\tightlist
\item
  \textbf{Label Propagation}: Propagates labels across a graph using
  similarity kernels (like KNN or RBF).
\item
  \textbf{Pseudo-Labeling}: Fits a model on labeled data, predicts
  labels for unlabeled data, adds highly confident predictions to the
  training set, and repeats.
\end{enumerate}

\subsubsection{Programming Guide}\label{programming-guide-17}

```python from sklearn.semi\_supervised import LabelPropagation import
numpy as np

\section{Unlabeled data points are marked with
-1}\label{unlabeled-data-points-are-marked-with--1}

y\_mixed = np.array({[}0, 1, -1, -1, 0, -1{]})

lp = LabelPropagation(kernel=`knn', n\_neighbors=2) lp.fit(X, y\_mixed)
complete\_labels = lp.transduction\_ ```\n\n---\n\n\# Chapter 19:
Classification on Imbalanced Data

Imbalanced data occurs when one class is extremely rare compared to the
other.

\subsubsection{High-Level Intuition}\label{high-level-intuition-18}

If only 0.1\% of bank transactions are fraudulent, a model that simply
outputs ``Not Fraud'' achieves 99.9\% accuracy but is useless. We must
balance the classes or adjust the model's cost function.

\subsubsection{Core Concepts}\label{core-concepts-11}

\begin{enumerate}
\def\labelenumi{\arabic{enumi}.}
\tightlist
\item
  \textbf{Imbalance Trap}: Accuracy is misleading for skewed datasets.
\item
  \textbf{Oversampling}: Duplicating minority class samples to balance
  the training set.
\item
  \textbf{Cost-Sensitive Learning}: Penalizing minority class errors
  more heavily during training.
\end{enumerate}

\subsubsection{Programming Guide}\label{programming-guide-18}

```python from sklearn.svm import SVC from sklearn.metrics import
f1\_score

\section{Using class\_weight=`balanced' to penalize rare class
errors}\label{using-class_weightbalanced-to-penalize-rare-class-errors}

model = SVC(class\_weight=`balanced', random\_state=42)
model.fit(X\_train, y\_train)

\section{F1-score evaluation}\label{f1-score-evaluation}

score = f1\_score(y\_test, model.predict(X\_test), average=`macro')
```\n\n---\n\n\# Chapter 20: Hyperparameter Optimization with Optuna

Tuning models manually is slow. In this unit, we explore automated
optimization strategies.

\subsubsection{High-Level Intuition}\label{high-level-intuition-19}

\textbf{Optuna} is a framework that automates the search for model
hyperparameters. Instead of guessing randomly, it uses \textbf{Bayesian
Optimization} - constructing a statistical model of the parameter space
to guess which settings are likely to yield the best results.

\subsubsection{Core Concepts}\label{core-concepts-12}

\begin{enumerate}
\def\labelenumi{\arabic{enumi}.}
\tightlist
\item
  \textbf{Bayesian Optimization}: Uses past trial results to build a
  probability model of the objective function.
\item
  \textbf{Gaussian Process}: A regression method used to estimate
  performance and uncertainty across parameters.
\end{enumerate}

\subsubsection{Programming Guide}\label{programming-guide-19}

```python import optuna from sklearn.ensemble import
RandomForestClassifier

def objective(trial): n\_est = trial.suggest\_int(`n\_estimators', 10,
50) depth = trial.suggest\_int(`max\_depth', 2, 5) model =
RandomForestClassifier(n\_estimators=n\_est, max\_depth=depth)
model.fit(X\_train, y\_train) return model.score(X\_test, y\_test)

study = optuna.create\_study(direction=`maximize')
study.optimize(objective, n\_trials=10) ```\n\n---\n\n\# Chapter 21:
Reinforcement Learning (RL)

Reinforcement Learning trains agents to make decisions by rewarding
desired behaviors and punishing undesired ones.

\subsubsection{High-Level Intuition}\label{high-level-intuition-20}

This is learning by trial-and-error (like training a dog). An agent
takes actions in an environment, receives rewards, and updates a memory
table (\textbf{Q-Table}) of action values.

\subsubsection{Core Concepts \& Math}\label{core-concepts-math-3}

\begin{enumerate}
\def\labelenumi{\arabic{enumi}.}
\tightlist
\item
  \textbf{Markov Decision Process (MDP)}: States, Actions, Rewards, and
  Transitions.
\item
  \textbf{Exploration vs.~Exploitation}: Balancing trying new actions
  (epsilon-greedy exploration) and taking the best known actions.
\item
  \textbf{Bellman Equation (Q-Update)}:
  \[Q(s, a) \leftarrow Q(s, a) + \alpha \left[ R + \gamma \max_{a'} Q(s', a') - Q(s, a) \right]\]
  Where \(\alpha\) is the learning rate, \(\gamma\) is the discount
  factor, and \(R\) is the reward.
\end{enumerate}

\subsubsection{Programming Guide}\label{programming-guide-20}

\texttt{python\ \#\ Q-learning\ update\ step\ td\_target\ =\ reward\ +\ gamma\ *\ np.max(q\_table{[}next\_state{]})\ q\_table{[}state,\ action{]}\ +=\ alpha\ *\ (td\_target\ -\ q\_table{[}state,\ action{]})}\n\n---\n\n\#
Chapter 22: Deep Computer Vision \& Convolutional Networks

Computer vision uses deep neural networks to extract information from
images.

\subsubsection{High-Level Intuition}\label{high-level-intuition-21}

Images are grids of pixels. To recognize shapes, we slide a small grid
of weights (a \textbf{Convolutional Filter}) across the image. This
filter highlights matching attributes (like edges), which are then
compressed using \textbf{Max Pooling} to keep the representation
compact.

\subsubsection{Core Concepts}\label{core-concepts-13}

\begin{enumerate}
\def\labelenumi{\arabic{enumi}.}
\tightlist
\item
  \textbf{2D Convolution}: Multiplying a filter matrix over pixel
  neighborhoods to construct feature maps.
\item
  \textbf{Max Pooling}: Sliding a window across a feature map and
  retaining only the maximum value, reducing image size.
\item
  \textbf{Sobel Filter}: A specific filter used to detect horizontal or
  vertical edges.
\end{enumerate}

\subsubsection{Programming Guide}\label{programming-guide-21}

```python \# 2D Convolution operation from scratch output{[}i, j{]} =
np.sum(image{[}i:i+f, j:j+f{]} * kernel)

\section{Max pooling operation}\label{max-pooling-operation}

output\_pool{[}i, j{]} = np.max(feature\_map{[}i\emph{2:i}2+2,
j\emph{2:j}2+2{]}) ```\n\n---\n\n\# Chapter 23: Graph Machine Learning
\& Link Analysis

Graph ML studies datasets structured as networks (nodes and edges).

\subsubsection{High-Level Intuition}\label{high-level-intuition-22}

Many real-world datasets are networks of interconnected items (like
social media friends or web pages). \textbf{Graph ML} studies these
relationships, predicting node roles or link connections.

\subsubsection{Core Concepts}\label{core-concepts-14}

\begin{enumerate}
\def\labelenumi{\arabic{enumi}.}
\tightlist
\item
  \textbf{Adjacency Matrix (\(A\))}: Grid showing node connections.
\item
  \textbf{Degree Matrix (\(D\))}: Diagonal grid showing connection
  counts.
\item
  \textbf{Graph Laplacian (\(L = D - A\))}: Used to study graph
  structures.
\item
  \textbf{PageRank}: Computes node importance based on link structures.
\end{enumerate}

\subsubsection{Programming Guide}\label{programming-guide-22}

\texttt{python\ \#\ Creating\ Laplacian\ matrix\ A\ =\ np.array({[}{[}0,\ 1,\ 0{]},\ {[}1,\ 0,\ 1{]},\ {[}0,\ 1,\ 0{]}{]})\ D\ =\ np.diag(np.sum(A,\ axis=1))\ L\ =\ D\ -\ A\ \ \#\ Graph\ Laplacian}\n\n---\n\n\#
Chapter 24: Model Interpretability \& Explainable AI (XAI)

XAI focuses on explaining the predictions of complex models.

\subsubsection{High-Level Intuition}\label{high-level-intuition-23}

If a model makes predictions, we need to understand why. *
\textbf{Permutation Importance} shuffles a feature's values. If accuracy
drops significantly, that feature is important. * \textbf{SHAP Values}
use cooperative game theory to measure the contribution of each feature
to a prediction.

\subsubsection{Programming Guide}\label{programming-guide-23}

```python from sklearn.inspection import permutation\_importance

\section{Extracting permutation
importance}\label{extracting-permutation-importance}

result = permutation\_importance(model, X\_val, y\_val, n\_repeats=5)
importances = result.importances\_mean ```\n\n---\n\n\# Project 1:
Predict Customer Churn (Classification Case Study)

In this project, you will build a binary classification pipeline to
predict customer churn using demographic and transaction data.

\subsubsection{Workflow Steps:}\label{workflow-steps}

\begin{enumerate}
\def\labelenumi{\arabic{enumi}.}
\tightlist
\item
  Load dataset and calculate base churn rate.
\item
  Encode categorical columns and scale numerical features.
\item
  Train a Random Forest classifier.
\item
  Tune probability thresholds using Precision-Recall curves to balance
  precision and recall.
\end{enumerate}

\subsubsection{Code Checklist:}\label{code-checklist}

\begin{itemize}
\tightlist
\item
  \texttt{df{[}\textquotesingle{}Churn\textquotesingle{}{]}.mean()}
\item
  \texttt{pd.get\_dummies()}
\item
  \texttt{MinMaxScaler()}
\item
  \texttt{RandomForestClassifier()}\n\n---\n\n\# Project 2: Housing
  Price Prediction (Regression Case Study)
\end{itemize}

In this project, you will build a regression pipeline to predict housing
prices.

\subsubsection{Workflow Steps:}\label{workflow-steps-1}

\begin{enumerate}
\def\labelenumi{\arabic{enumi}.}
\tightlist
\item
  Load dataset and handle missing values.
\item
  Construct polynomial features to capture non-linear relationships.
\item
  Fit regularized Ridge and Lasso regression models.
\item
  Optimize hyperparameters using Optuna.
\end{enumerate}

\subsubsection{Code Checklist:}\label{code-checklist-1}

\begin{itemize}
\tightlist
\item
  \texttt{PolynomialFeatures()}
\item
  \texttt{Ridge(alpha=...)}
\item
  \texttt{optuna.create\_study()}\n\n---\n\n\# Project 3: Sentiment
  Analysis on Customer Reviews (NLP Case Study)
\end{itemize}

In this project, you will build a text classifier to analyze review
sentiment.

\subsubsection{Workflow Steps:}\label{workflow-steps-2}

\begin{enumerate}
\def\labelenumi{\arabic{enumi}.}
\tightlist
\item
  Preprocess reviews (remove punctuation, lowercase, strip stopwords).
\item
  Vectorize text with a TF-IDF transformer.
\item
  Train a Multinomial Naive Bayes model.
\item
  Run permutation feature importance to identify key sentiment words.
\end{enumerate}

\subsubsection{Code Checklist:}\label{code-checklist-2}

\begin{itemize}
\tightlist
\item
  \texttt{TfidfVectorizer()}
\item
  \texttt{MultinomialNB()}
\item
  \texttt{permutation\_importance()}\n\n---\n\n\# Project 4: Iris
  Species Classification (Multi-class Case Study)
\end{itemize}

In this project, you will build a multi-class classification pipeline to
predict flower species.

\subsubsection{Workflow Steps:}\label{workflow-steps-3}

\begin{enumerate}
\def\labelenumi{\arabic{enumi}.}
\tightlist
\item
  Load the dataset and calculate species frequencies.
\item
  Map target text labels to integers.
\item
  Apply standard scaling to numeric features.
\item
  Train and compare Decision Tree and Naive Bayes classifiers using
  class-specific F1-scores.
\end{enumerate}

\subsubsection{Code Checklist:}\label{code-checklist-3}

\begin{itemize}
\tightlist
\item
  \texttt{df{[}\textquotesingle{}Species\textquotesingle{}{]}.map()}
\item
  \texttt{StandardScaler()}
\item
  \texttt{DecisionTreeClassifier()}
\item
  \texttt{GaussianNB()}\n\n---\n\n\# Project 5: Customer Segmentation
  (Unsupervised Case Study)
\end{itemize}

In this project, you will segment customers into distinct profiles using
clustering.

\subsubsection{Workflow Steps:}\label{workflow-steps-4}

\begin{enumerate}
\def\labelenumi{\arabic{enumi}.}
\tightlist
\item
  Load dataset and scale demographic features.
\item
  Apply the Elbow Method to choose \(K\).
\item
  Run K-Means clustering.
\item
  Map clusters to a 2D plane using PCA and profile customer groups based
  on cluster means.
\end{enumerate}

\subsubsection{Code Checklist:}\label{code-checklist-4}

\begin{itemize}
\tightlist
\item
  \texttt{MinMaxScaler()}
\item
  \texttt{KMeans(n\_clusters=...)}
\item
  \texttt{PCA(n\_components=2)}
\item
  \texttt{df.groupby(\textquotesingle{}Cluster\textquotesingle{}).mean()}\n\n---\n\n\#
  Project 6: Daily Sales Forecasting (Time Series Case Study)
\end{itemize}

In this project, you will forecast daily sales using past transaction
records.

\subsubsection{Workflow Steps:}\label{workflow-steps-5}

\begin{enumerate}
\def\labelenumi{\arabic{enumi}.}
\tightlist
\item
  Load time series data and compute rolling means.
\item
  Check stationarity using the ADF test.
\item
  Apply differencing to remove trends.
\item
  Fit an ARIMA model and evaluate predictions using Mean Absolute
  Percentage Error (MAPE).
\end{enumerate}

\subsubsection{Code Checklist:}\label{code-checklist-5}

\begin{itemize}
\tightlist
\item
  \texttt{df{[}\textquotesingle{}Sales\textquotesingle{}{]}.rolling(window=7).mean()}
\item
  \texttt{adfuller()}
\item
  \texttt{df{[}\textquotesingle{}Sales\textquotesingle{}{]}.diff()}
\item
  \texttt{ARIMA(order=(1,\ 1,\ 1))}\n\n---\n\n\# Project 7: Credit Card
  Fraud Detection (Imbalanced Case Study)
\end{itemize}

In this project, you will detect rare fraudulent transactions.

\subsubsection{Workflow Steps:}\label{workflow-steps-6}

\begin{enumerate}
\def\labelenumi{\arabic{enumi}.}
\tightlist
\item
  Calculate the fraud ratio.
\item
  Balance the training set using random oversampling.
\item
  Train an SVM with balanced class weights.
\item
  Evaluate predictions using Precision-Recall Area Under the Curve
  (PR-AUC) and optimize decision thresholds.
\end{enumerate}

\subsubsection{Code Checklist:}\label{code-checklist-6}

\begin{itemize}
\tightlist
\item
  \texttt{random\_oversample()}
\item
  \texttt{SVC(class\_weight=\textquotesingle{}balanced\textquotesingle{})}
\item
  \texttt{precision\_recall\_curve()}
\item
  \texttt{auc()}\n\n---\n\n
\end{itemize}
